%%%%%%%%%%%%%%%%%%%%%%%%%%%%%%%%%%%%%%%%%

%
% Important note:
% This template requires the resume.cls file to be in the same directory as the
% .tex file. The resume.cls file provides the resume style used for structuring the
% document.

%
% Creator Peilin Li
% Contact me via twitter/wechat: @pe1l1nl1
% linkedin.com/peill and/or github/ppeill
% Inspired by Peppa Pig 
%%%%%%%%%%%%%%%%%%%%%%%%%%%%%%%%%%%%%%%%

%----------------------------------------------------------------------------------------
%	PACKAGES AND OTHER DOCUMENT CONFIGURATIONS
%----------------------------------------------------------------------------------------

\documentclass{cv} % Use the custom resume.cls style

\usepackage[left=0.40in,top=0.3in,right=0.75in,bottom=0.1in]{geometry} % Document margins
\usepackage{fontawesome}
\usepackage{times}
\usepackage{url}
%----------------------------------------------------------------------------------------
%	BIBLIOGRAPHY (publications.bib) -- compact style
%----------------------------------------------------------------------------------------
\usepackage[backend=bibtex, style=numeric, sorting=ydnt, giveninits=true, maxbibnames=99,
            doi=false, isbn=false, url=false, eprint=false]{biblatex}
\addbibresource{publications.bib}
\renewbibmacro{in:}{} % Drop "In:"
\AtEveryBibitem{\clearlist{publisher}\clearlist{location}\clearlist{organization}%
  \clearfield{pages}\clearfield{month}\clearfield{volume}\clearfield{number}\clearname{editor}}
\usepackage{csquotes}
\DefineBibliographyExtras{english}{\uspunctuation} % Commas/periods inside quotes
\renewcommand*{\newunitpunct}{\addcomma\space} % Comma-separated: Authors, "Title," Venue, Year.
\DeclareFieldFormat{journaltitle}{\mkbibemph{#1}}
\renewbibmacro*{issue+date}{\setunit{\addcomma\space}\usebibmacro{date}\newunit} % Article year without parens
\renewcommand*{\finalnamedelim}{\addspace\&\space} % "A, B \& C"
\renewcommand*{\mkbibcompletename}[1]{\ifdefstring{\namepartfamily}{Kurlandski}{\textbf{#1}}{#1}} % Bold my name
\defbibenvironment{bibliography} % Dash-bulleted hanging paragraphs matching the rest of the CV, no numbers
  {\setlength{\parskip}{1pt}\advance\leftskip 0.8em\vspace{3pt}}
  {\endgraf}
  {\endgraf\leavevmode\hspace*{-0.8em}\makebox[0.8em][l]{-}}
\usepackage{hyperref}
\newcommand{\tab}[1]{\hspace{.2667\textwidth}\rlap{#1}}
\newcommand{\itab}[1]{\hspace{0em}\rlap{#1}}

\name{Luke Kurlandski} % Your name 

\address{
\faEnvelope{\hspace{0.1ex} lk3591@g.rit.edu} \hspace{0.75ex}
\faGithub{\hspace{0.1ex} @lkurlandski}  \hspace{0.75ex}
\faLinkedin{\hspace{0.1ex} @luke-kurlandski} \hspace{0.75ex}
\faGraduationCap{\hspace{0.1ex} @Luke Kurlandski}
}

\begin{document}

%----------------------------------------------------------------------------------------
%	Education
%----------------------------------------------------------------------------------------
\begin{rSection}{Education}
	
{\bf Doctor of Philosophy (Ph.D.) @Rochester Institute of Technology (RIT) } \hfill {\em Aug. 2022 - current} 
\\ - Department: Computing and Information Science
\\ - Advised by: Dr. Matthew Wright \& Dr. Yin Pan

{\bf Bachelor of Science (B.S.) @The College of New Jersey (TCNJ) } \hfill {\em Aug. 2018 - May 2022} 
\\ - Department: Computer Science
\\ - Minors: Mathematics, Physics

\end{rSection}

%--------------------------------------------------------------------------------
%    Skills
%--------------------------------------------------------------------------------
\begin{rSection}{Skills}

\begin{tabular}{ @{} >{\bfseries}l @{\hspace{6ex}} l }
Languages: \ & Python, MatLab, R, Mathematica, Java, C++, C, Rust
\\ Libraries: \ & PyTorch, huggingface, scikit-learn, TensorFlow
\\ Tools: \ & Git, SLURM, GNU-Linux, LaTeX
\\ Security: \ & IDA, redare2, ghidra, OllyDbg
\end{tabular}

\end{rSection}


%--------------------------------------------------------------------------------
%    Objective
%--------------------------------------------------------------------------------
\iffalse
\begin{rSection}{Objective}
To work for an organization which provides me the opportunity to improve my skills and knowledge to grow along with the organization objective.
\end{rSection}
\fi

%--------------------------------------------------------------------------------
%    Research
%--------------------------------------------------------------------------------
\begin{rSection}{Research Experience}
	
{\bf Malware \& ML, RIT}  \hfill {\em Aug. 2022 - current} 
\\{\textit{Graduate Research Assistant}; advised by Dr. Yin Pan \& Dr. Matthew Wright.}
\\- Researched malware classification and adversarial malware generation using state-of-the-art deep learning tech.

{\bf Learning Curve, TCNJ} \hfill {\em Feb. 2021 - May 2022} 
\\{\textit{Undergraduate Research Assistant}; advised by Dr. Michael Bloodgood.}
\\- Researched active learning stopping methods for low-resource natural language processing tasks.

{\bf Cybersecurity Division, CACI International} \hfill {\em Jun. 2021 - Aug. 2021} 
\\{\textit{Research Intern}}
\\- Researched deepfake audio detection for US Department of Defense (DoD) cybersecurity applications.

{\bf Collective Intelligence, TCNJ} \hfill {\em Jun. 2020 - Feb. 2021} 
\\{\textit{Undergraduate Research Assistant}; advised by Dr. Monisha Pulimood}
\\- Researched collective intelligence for web-applications and built machine learning recommendation system.

{\bf Photonics \& Materials Science, TCNJ} \hfill {\em Dec. 2018 - Feb. 2020} 
\\{\textit{Undergraduate Research Assistant}; advised by Dr. David McGee}
\\- Researched optical microstructures; wrote APIs for laboratory hardware and GUI applications for experiments.
\end{rSection}

%--------------------------------------------------------------------------------
%    Publications
%--------------------------------------------------------------------------------
%\iffalse

\begin{rSection}{Publications}
\nocite{*}
\printbibliography[heading=none]
\end{rSection}

\pagebreak

\begin{rSection}{Posters \& Talks}
	
	- \textbf{Kurlandski, L.}, Pan, Y. Wright, M. ``From Bytes to Insights: A New Era of Malware Detection''. In preparation to be presented at \textit{Western New York Image and Signal Processing Workshop (WNYISPW)}. 2023.
	
	- \textbf{Kurlandski, L.} ``Transferring Benignness to Assembly Instructions''. Presented in fulfillment of RIT's research potential assessment for Ph.D. students. 2023.
	
	- \textbf{Kurlandski, L.}, Pan, Y. Wright, M. ``An Adversarial Malware Evasion Attack Using Machine Translation''. Presented at \textit{AI@RIT Summit}. 2022.

	- \textbf{Kurlandski, L.} ``Detecting Deepfake Audio with Deep Neural Networks''. Presented in fulfillment of TCNJ's student capstone. 2022.

\end{rSection}

%--------------------------------------------------------------------------------
%    Projects
%--------------------------------------------------------------------------------
%\iffalse
\begin{rSection}{Projects}
	
	{\textbf{MDALh}} \hfill {\em Jun. 2023 - current} 
	\\{\small \textit{\url{github.com/lkurlandski/MDALTH}}}
	\\- A modular deep active learning library with high-level wrappers for huggingface pretrained models.
	\\ - Used Python with PyTorch, transformers, datasets, tokenizers, etc.

	{\bf Explaining MalConv} \hfill  {\em Oct. 2022 - current} 
    \\{\small \textit{\url{github.com/lkurlandski/MalConv2}}}
	\\- Identify notable malware byte patterns for MalConv malware classifier using explainable AI techniques.
	\\- Used Python, C++ with PyTorch, captum, pefile, LIEF, capstone, etc.

	{\bf Adversarial Evasion Attacks}  \hfill  {\em Apr. 2023} 
	\\{\small \textit{\url{github.com/lkurlandski/AdversarialEvasion}}}
	\\- Implemented FGSM and PGD adversarial attacks to evade pretrained image classifiers.
	\\- Used Python with PyTorch, etc.

	{\bf Website Traffic Fingerprinting}  \hfill  {\em Mar. 2023} 
	\\{\small \textit{\url{github.com/lkurlandski/WebsiteFingerprinting}}}
	\\- Trained deep learning classifiers to perform website traffic fingerprinting.
	\\- Used Python with PyTorch, TensorFlow, etc.

	{\bf scikit-AL}  \hfill  {\em Jan. 2022 - Aug. 2022} 
    \\{\small \textit{\url{github.com/lkurlandski/ActiveLearning}}}
	\\- A framework for analyzing active learning stopping methods and query algorithms.
	\\- Used Python with scikit-learn, huggingface, ModAL, etc.

	{\bf Collective Recommender}  \hfill  {\em May 2020 - Nov. 2020} 
	\\- A hybrid recommendation system to help users find relevant academic literature.
	\\- Used Python \& PostgresSQL with LightFM, NetworkX, pandas, psycopg2, etc.

   	{\bf Tchisla Solver}  \hfill  {\em Jan. 2020 - May 2020} 
    \\{\small \textit{\url{https://github.com/lkurlandski/TchislaSolver}}}
    \\{\small \textit{\url{https://luke-kurlandski.shinyapps.io/TchislaSolver/}}}
    \\- A dynamic programming algorithm and web application to find solutions for the math puzzle game, \textit{Tchisla}.
    \\- Used Python \& R

	{\bf Newport Hologram Creation}  \hfill  {\em Jan. 2019 - May 2019} 
    \\{\small \textit{\url{github.com/lkurlandski/Newport-Hologram-Creation}}}
	\\- An API to control laboratory equipment and a GUI application to run experiments that create holograms.
	\\- Used Python with PySerial, tkinter, etc.

\end{rSection}

\pagebreak

\begin{rSection}{Teaching \& Leadership}

	{\textbf{Instructor, RIT}} \hfill {\em Jan. 2024 - May 2024} 
    \\- Instructed graduate research seminar on Advanced Malware Forensics.

	{\textbf{Graduate Teaching Assistant, RIT}} \hfill {\em Aug. 2023 - Dec. 2024} 
    \\- Aid for graduate research seminar on Human Factors in Security and Privacy.

   	{\textbf{Research Mentor, RIT}} \hfill {\em Jan. 2023 - current} 
    \\- Mentored undergraduate students for capstone projects, summer research experiences, and independent studies.

   	{\textbf{Lead Undergraduate Research Assistant, TCNJ}} \hfill {\em Aug. 2021 - May. 2022} 
    \\- Student lead among undergraduate RAs (Learning Curve Group).

   	{\textbf{GED Tutor, TCNJ}} \hfill {\em Nov. 2017 - May. 2018} 
    \\- Taught math/science/english to adults seeking to earn their general equivalency degree.

\end{rSection}

\begin{rSection}{Awards \& Honors}

    {\textbf{National Science Foundation (NSF) Graduate Research Fellowship Program (GRFP) (recipient)}} \hfill {\em Apr. 2024} 
    \\ - Recipient of prestigious national fellowship program.

    {\textbf{National Defense Science and Engineering Graduate (NDSEG) Fellowship (recipient)}} \hfill {\em Mar. 2024} 
    \\ - Recipient of prestigious national fellowship program. Declined due to conflict with GRFP.

    {\textbf{AI Seed Fund (recipient), RIT}} \hfill {\em Jan. 2024} 
    \\ - Recipient of local funding competition at RIT.

    {\textbf{Goldberg-Neff Scholarship (recipient), TCNJ}} \hfill {\em May 2022} 
    \\ - Recipient of scholarship for graduating students pursuing graduate research.

	{\textbf{Upsilon Pi Epsilon Computer Science Honor's Society (member), TCNJ}} \hfill {\em May 2022}
    \\ - Inducted as a member into the distinguished national Honor's Society.

	{\textbf{Phi Beta Kappa Honor's Society (member), TCNJ}} \hfill {\em May 2022} 
    \\ - Inducted as a member into the distinguished national Honor's Society.

	{\textbf{Kappa Delta Rho Nu Gamma Chapter (founder), TCNJ}} \hfill {\em Jan. 2019 - May 2022} 
    \\ - Established the Nu Gamma Chapter of Kappa Delta Rho fraternity at TCNJ.

	{\textbf{Department of Computer Science Award (recipient), TCNJ}} \hfill {\em May 2019}
    \\ - Recipient of department award for first-year students.

	{\textbf{Dean's List 8/8 semesters, TCNJ}} \hfill {\em Aug. 2018 - May 2022}
    \\ - Maintained a GPA above 3.5 all eight semesters of undergraduate.

\end{rSection}

\end{document}